\documentclass[11pt]{article}
\usepackage[margin=1in]{geometry}
\usepackage{times}
\usepackage{booktabs}
\usepackage{graphicx}
\usepackage{amsmath}
\usepackage{hyperref}
\usepackage{xcolor}
\usepackage{listings}
\usepackage{microtype}
\usepackage{natbib}
\usepackage{pgfplotstable}
\usepackage{pgfplots}
\pgfplotsset{compat=1.18}

% JSON listings style
\lstdefinelanguage{json}{
  basicstyle=\ttfamily\small,
  morestring=[b]",
  stringstyle=\color{teal},
  literate=
    {true}{{\color{violet}true}}{4}
    {false}{{\color{violet}false}}{5},
}

\hypersetup{colorlinks=true, linkcolor=blue, citecolor=blue, urlcolor=blue}

\title{\textbf{Format Anchoring Bias in LLM-as-Judge:\\
How a Single Token in the Output Format\\
Corrupts Automated Code Evaluation}}

\author{
  Srinivas Vaddi\\
  \texttt{vaddisrinivas@gmail.com}
}

\date{February 2026}

\begin{document}

\maketitle

\begin{abstract}
LLM-as-judge is widely used to evaluate code quality, but we identify a
previously-unreported failure mode: \emph{format anchoring bias}.
When a judge prompt includes a concrete boolean value in the output format
example --- such as \texttt{"verdict":true} --- small language models copy
that example value in up to \textbf{91.3\%} of responses regardless of the
code's actual correctness, collapsing evaluation accuracy to \textbf{18.2\%}.
We demonstrate this on a 242-run evaluation suite (22 bug-fixing tasks,
11 prompt variants, 1 model) using pytest execution as ground truth.
Fixing three prompt anti-patterns --- (1) anchored format examples,
(2) ground-truth leakage via an \texttt{expected} field,
and (3) negative-question framing --- recovers accuracy to \textbf{92.1\%}
with the same judge model (GPT-5-mini), a \textbf{+69.4 percentage-point}
improvement with zero model upgrade.
We further identify a \emph{reasoning-verdict dissociation} effect:
in 48 cases the judge's chain-of-thought explicitly states the code is
correctly fixed yet the verdict JSON token is still \texttt{true},
suggesting the format anchor operates \emph{after} the reasoning trace,
at JSON generation time.
All data and code are publicly available.
\end{abstract}

%-----------------------------------------------------------------------------
\section{Introduction}
%-----------------------------------------------------------------------------

LLM-as-judge~\citep{zheng2023judging} has become the dominant paradigm for
automated evaluation of code generation and bug-fixing systems, offering
scalability and flexibility that static test suites cannot match.
However, its reliability depends entirely on the quality of the judge
prompt --- a dependency that is poorly understood in practice.

We report a simple but severe failure mode that arises when the judge
prompt includes a concrete boolean value in the structured output format
specification.
For example, a prompt that instructs the judge to return
\texttt{\{"verdict":true,...\}} as an example output causes small models
to produce \texttt{true} in the vast majority of responses, independent
of the code under evaluation.
We call this \emph{format anchoring bias}.

The practical impact is large.
In our evaluation of a code bug-fixing benchmark (242 runs, pytest oracle),
a biased judge prompt produced only 18.2\% accuracy against ground truth;
fixing the prompt --- without changing the judge model --- recovered 90.1\%
accuracy.
Upgrading to a stronger judge model (Claude Sonnet 4.6) yielded a further
modest gain to 92.1\%, but the prompt fix accounted for the overwhelming
majority of the improvement.

Our contributions are:
\begin{itemize}
  \item We identify and characterise \emph{format anchoring bias} in
        LLM-as-judge systems for code evaluation (\S\ref{sec:bias}).
  \item We provide a controlled 4-condition ablation isolating the
        prompt effect from the model effect (\S\ref{sec:results}).
  \item We report \emph{reasoning-verdict dissociation}: the judge's
        chain-of-thought reaches the correct conclusion but the emitted
        JSON token contradicts it (\S\ref{sec:dissociation}).
  \item We release a complete reproducibility package: 242 evaluation runs
        with four scoring conditions, all judge payloads, pytest outputs,
        and analysis scripts.
\end{itemize}

%-----------------------------------------------------------------------------
\section{Background}
%-----------------------------------------------------------------------------

\paragraph{LLM-as-judge.}
\citet{zheng2023judging} introduced the paradigm of using a strong LLM
to score model outputs, showing high correlation with human judgments.
Subsequent work has identified various biases including position
bias~\citep{wang2023large}, verbosity bias~\citep{saito2023verbosity},
and self-enhancement bias~\citep{panickssery2024llm}.
Format anchoring bias is distinct: it arises not from the content of the
response being judged, but from a token in the \emph{judge prompt itself}.

\paragraph{Structured output generation.}
When LLMs are asked to produce JSON, they are sensitive to format
examples~\citep{brown2020language}.
In few-shot prompting, models copy the statistical patterns of examples;
a JSON format example with a specific value acts as a strong prior on
the corresponding position in the output.

\paragraph{Code evaluation benchmarks.}
SWE-bench~\citep{jimenez2024swebench} and its variants evaluate code fixes
against repository test suites.
Our setting differs: we use a controlled benchmark (VoltSnip BUG01--BUG22)
with hand-authored constraints, enabling us to separate the judge's
behaviour from the test suite's coverage.

%-----------------------------------------------------------------------------
\section{The VoltSnip Evaluation Framework}
\label{sec:framework}
%-----------------------------------------------------------------------------

\paragraph{Tasks.}
We use 22 code bug-fixing tasks (BUG01--BUG22) drawn from a single
FastAPI service (\texttt{orgops}).
Each task specifies a target file, line range, a buggy code snippet, and
a natural-language description of the fix required.
Task categories include \textit{http\_resilience},
\textit{error\_contracts}, \textit{db\_patterns},
\textit{logging\_privacy}, and \textit{concurrency\_cache}.

\paragraph{Variants.}
We test 11 prompt variants (P0--P10) that vary: presence of oracle criteria
in the prompt, whether retrieved snippets are injected as context, and
whether the model has access to retrieval tools.
This gives $22 \times 11 = 242$ runs.

\paragraph{Model.}
All runs use \texttt{claude-haiku-4-5} via the ClaudeCode provider.
The same model and code outputs are rescored across all four judge
conditions; the model never re-runs.

\paragraph{Ground truth.}
Each task includes a pytest test command
(e.g., \texttt{pytest -k bug\_09}).
The test result (pass/fail) is our ground truth oracle, executed in a
Docker container against the model's patched file.
Pytest pass rate: \textbf{96.3\%} (233/242).

\paragraph{Constraint-based scoring.}
Each task defines one or more constraints, each with an \texttt{id},
\texttt{check} description, \texttt{judge\_prompt}, and
\texttt{expected} boolean.
\texttt{expected=false} constraints check that a failure mode is
\emph{absent} (the bug is fixed); \texttt{expected=true} constraints
check that a positive property is \emph{present}.
A run passes if the fraction of satisfied constraints meets a threshold
($\tau = 0.7$).

%-----------------------------------------------------------------------------
\section{Format Anchoring Bias}
\label{sec:bias}
%-----------------------------------------------------------------------------

\subsection{Three Prompt Anti-Patterns}

Our original judge prompt contained three compounding errors:

\paragraph{Anti-pattern 1: Anchored format example.}
The output format specification included a concrete \texttt{true} value:
\begin{lstlisting}[language=json]
Return ONLY JSON:
{"results":[{"verdict":true,"reason":"..."},...]}
\end{lstlisting}
Small models treat format examples as strong priors.
GPT-4o-nano emitted \texttt{"verdict":true} in \textbf{91.3\%} of all
responses.

\paragraph{Anti-pattern 2: Ground-truth leakage.}
Each constraint was sent to the judge including its \texttt{expected}
field:
\begin{lstlisting}[language=json]
{"id": "bug09_failure_absent",
 "check": "Status mapping remains incorrect.",
 "expected": false,   // <-- leaked!
 "judge_prompt": "Does the implementation still exhibit..."}
\end{lstlisting}
For \texttt{expected=false} constraints, a \texttt{true} verdict means
the bug is \emph{still present} and the run \emph{fails}.
Combined with Anti-pattern 1, this made the bias maximally destructive
for the most common constraint type.

\paragraph{Anti-pattern 3: Negative-question framing.}
The \texttt{judge\_prompt} was phrased as:
\textit{``Does the implementation \textbf{still} exhibit this failure
mode?''}, priming the model toward \texttt{true} (``yes, it does'').
A positive framing (``Answer true if the described behavior IS present'')
is less susceptible.

\subsection{Why Format Anchors Dominate}

Few-shot examples in LLM prompts establish strong statistical priors on
output tokens~\citep{brown2020language}.
When the output is a structured JSON object, the model must generate each
field value conditioned on the prior context.
A concrete boolean in the example (\texttt{true}) occupies the same
token position as the verdict the model must emit.
For small models with limited instruction-following capacity, the format
example overwhelms the semantic signal from the code and constraint.

%-----------------------------------------------------------------------------
\section{Experimental Setup}
%-----------------------------------------------------------------------------

We compare four scoring conditions on the \emph{same} 242 runs:

\begin{enumerate}
  \item \textbf{GPT-4o-nano (biased)}: original prompt with all three
        anti-patterns. Judge: \texttt{gpt-4o-nano} (the original
        production judge).
  \item \textbf{GPT-5-mini (biased)}: same biased prompt, stronger judge.
  \item \textbf{GPT-5-mini (fixed)}: fixed prompt, same model as (2).
        Isolates prompt effect.
  \item \textbf{Claude Sonnet 4.6 (fixed)}: fixed prompt, strongest judge.
        Provides a ceiling estimate.
\end{enumerate}

The fixed prompt (conditions 3--4):
\begin{lstlisting}[language=json, breaklines=true]
"Judge each constraint independently. Answer true if the
described behavior IS present; false if NOT present.
Return ONLY JSON:
{"results":[{"verdict":true,"reason":"..."},
            {"verdict":false,"reason":"..."},...]}
verdict must be a JSON boolean, not a string."
\end{lstlisting}
Key changes: (a) format example shows \emph{both} \texttt{true} and
\texttt{false}; (b) \texttt{expected} field removed from constraint
payloads; (c) affirmative framing replaces negative-question framing.

%-----------------------------------------------------------------------------
\section{Results}
\label{sec:results}
%-----------------------------------------------------------------------------

\begin{table}[t]
\centering
\caption{
  Scoring accuracy across four judge conditions.
  All conditions evaluate the same 242 code outputs.
  Ground truth: pytest oracle (96.3\% pass rate).
  vTrue\% = fraction of judge verdicts that were \texttt{true}.
  Acc\% = agreement with pytest oracle.
}
\label{tab:main}
\begin{tabular}{lrrrrrrr}
\toprule
Condition & Pass\% & vTrue\% & TP & TN & FP & FN & Acc\% \\
\midrule
GPT-4o-nano (biased)      & 17.8 & 91.3 &  39 &  5 &  4 & 194 & 18.2 \\
GPT-5-mini (biased)       & 20.2 & 88.8 &  45 &  5 &  4 & 188 & 20.7 \\
\midrule
GPT-5-mini (fixed)        & 92.1 & 16.9 & 216 &  2 &  7 &  17 & 90.1 \\
Claude Sonnet 4.6 (fixed) & 93.4 & 15.7 & 220 &  3 &  6 &  13 & 92.1 \\
\midrule
Pytest oracle             & 96.3 &   —  & —   &  — &  — &   — & 100.0 \\
\bottomrule
\end{tabular}
\end{table}

Table~\ref{tab:main} reports the main results.
With the biased prompt, both GPT-4o-nano and GPT-5-mini produce
\texttt{true} verdicts over 88\% of the time and achieve less than 21\%
accuracy.
Fixing the prompt --- with the \emph{same} GPT-5-mini model --- recovers
90.1\% accuracy (+\textbf{69.4 pp}).
Upgrading to Claude Sonnet 4.6 adds a further 2.0 pp.

\paragraph{Statistical significance.}
McNemar's test~\citep{mcnemar1947note} on GPT-4o-nano (biased) vs.\
Claude Sonnet 4.6 (fixed):
$\chi^2(1) = 164.8$, $p \ll 0.001$ ($b=9$, $c=192$).
Cohen's $g = 0.864$ on verdict=\texttt{true} rates
(91.3\% vs.\ 15.7\%), indicating a very large effect.

\paragraph{Prompt effect dominates model effect.}
Comparing conditions 2 and 3 (same model, different prompt):
accuracy increases from 20.7\% to 90.1\% (+69.4 pp).
Comparing conditions 3 and 4 (same fixed prompt, model upgrade):
accuracy increases from 90.1\% to 92.1\% (+2.0 pp).
\emph{The prompt fix accounts for 97\% of the total recovery.}

%-----------------------------------------------------------------------------
\section{Reasoning-Verdict Dissociation}
\label{sec:dissociation}
%-----------------------------------------------------------------------------

A striking secondary finding is \emph{reasoning-verdict dissociation}:
cases where the judge's chain-of-thought explicitly concludes the code
is correctly fixed, yet the emitted verdict is still \texttt{true}
(``bug present'').

In the biased GPT-5-mini condition, we identify \textbf{48 such cases}
by matching judge reasons containing words such as
\textit{``correctly fixed'', ``no longer exhibits'', ``does not exhibit'',
``now maps properly''} against the emitted \texttt{"verdict":true}.

Example (BUG09, P0):
\begin{lstlisting}[language=json]
{
  "verdict": true,
  "reason": "The code now maps TimeoutError to 504 and
    ValueError to 400, fixing the status mapping."
}
\end{lstlisting}
The reason is entirely correct; the pytest test passes.
Yet the verdict contradicts the reason.

This pattern suggests that format anchoring operates at the \emph{token
generation} level, \emph{after} the chain-of-thought is complete.
The model reasons correctly in prose, then ignores its own conclusion when
generating the structured output field.
This is consistent with findings that chain-of-thought does not always
constrain subsequent token predictions in small
models~\citep{turpin2023language}.

After the prompt fix, dissociation cases drop to 2--3 across all 242
runs, confirming that the anchor was the cause.

%-----------------------------------------------------------------------------
\section{Implications for Practice}
%-----------------------------------------------------------------------------

\paragraph{Audit structured-output prompts for concrete example values.}
Any boolean, integer, or enum in a format example acts as a potential
anchor.
Show multiple distinct values (\texttt{true}/\texttt{false}) or use
abstract placeholders (\texttt{\textless{}bool\textgreater{}}).

\paragraph{Never leak expected values to the judge.}
Fields such as \texttt{expected}, \texttt{ground\_truth}, or
\texttt{correct\_answer} in the judge payload leak the target answer.
Even if the judge ``should'' ignore them, small models may not.

\paragraph{Use affirmative framing.}
Negative questions (``Does it \emph{still} exhibit...?'') prime
\texttt{true} responses.
Affirmative framing (``Is this behavior present?'') is more neutral.

\paragraph{Validate judge accuracy against an oracle.}
Pytest (or any deterministic oracle) provides cheap ground-truth
validation.
A judge achieving less than 50\% agreement with a 96\%+ passing oracle
is almost certainly biased, not measuring code quality.

%-----------------------------------------------------------------------------
\section{Limitations}
%-----------------------------------------------------------------------------

Our study uses a single model family for code generation
(\texttt{claude-haiku-4-5}) and a single codebase (\texttt{orgops}).
Whether format anchoring bias generalises to other code domains, other
generation models, or multi-constraint evaluations with heterogeneous
types requires further investigation.
Our contradiction detection heuristic (keyword matching on judge reasons)
likely undercounts dissociation cases; a manual audit on a random sample
would provide a tighter estimate.

%-----------------------------------------------------------------------------
\section{Conclusion}
%-----------------------------------------------------------------------------

We have demonstrated that a single concrete boolean token in a judge
prompt's output format specification can reduce LLM-as-judge accuracy
from 96\% (pytest oracle) to 18\% --- a catastrophic degradation that
is invisible without an independent oracle.
The fix is three simple prompt changes requiring no model upgrade and no
retraining.
We recommend that all LLM-as-judge pipelines using structured output be
audited for format anchoring before deployment, and that oracle
validation (even partial) be treated as a mandatory component of
evaluation harness design.

%-----------------------------------------------------------------------------
\bibliographystyle{abbrvnat}
\bibliography{references}

\appendix
\section{Reproducibility}
\label{app:repro}

All materials are available at:\\
\url{https://github.com/vaddisrinivas/voltsnip}

\begin{itemize}
  \item Branch \texttt{vsevals}: eval framework source
        (\texttt{vsevals/vsevals/scorer.py}, task YAML files,
        rescore scripts).
  \item Branch \texttt{paper-judge-failure-0-pytest}: 242 run artifacts,
        four \texttt{full\_dump\_\_\{suffix\}.json} per run,
        four \texttt{matrix\_results\_\_\*.csv} aggregates.
\end{itemize}

To reproduce the rescoring:
\begin{lstlisting}[basicstyle=\ttfamily\small]
git clone https://github.com/vaddisrinivas/voltsnip
cd voltsnip && git checkout vsevals
pip install -e vsevals/
python vsevals/scripts/rescore_scoring.py \
  --matrix-dir vsevals_runs/matrix_20260228T053125139679Z \
  --suite vsevals/suite.yaml \
  --judge-model openai:gpt-5-mini \
  --output-suffix my-rescore \
  --workers 4
\end{lstlisting}

To reproduce Table~\ref{tab:main}:
\begin{lstlisting}[basicstyle=\ttfamily\small]
git checkout paper-judge-failure-0-pytest
python paper/compute_stats.py
\end{lstlisting}

\end{document}
